\documentclass[12pt,a4paper]{article}
\usepackage[margin=1in]{geometry}
\usepackage{amsmath,amssymb,amsfonts}
\usepackage{algorithm}
\usepackage{algpseudocode}
\usepackage{enumitem}
\usepackage{tikz}
\usepackage{array}

\title{\textbf{CSE 317: Design and Analysis of Algorithms} \\ Homework 2}
\author{Abdullah Khan Sherwani \\ ERP: 27880}
\date{}

\begin{document}
\maketitle

\section*{Question 1: Triangle Maximum Path Sum}

\subsection*{Problem Statement}
Consider a triangular array of integers with $n$ rows, where the $i$-th row contains exactly $i$ positive numbers. Starting from the top element, at each step from position $j$ in row $i$, we may move to position $j$ or $j+1$ in row $i+1$. The goal is to find a path from top to bottom that maximizes the sum of selected numbers.

\subsection*{Dynamic Programming Formulation}

\subsubsection*{Notation}
Let $T[i][j]$ denote the value at position $j$ in row $i$ of the triangle, where $1 \le i \le n$ and $1 \le j \le i$.

\subsubsection*{Subproblem Definition}
Define $\text{dp}[i][j]$ as the \textbf{maximum path sum} starting from the top of the triangle $(T[1][1])$ and ending at position $j$ in row $i$.

\subsubsection*{Base Case}
\[
    \text{dp}[1][1] = T[1][1]
\]
There is only one element in the first row, so the maximum sum to reach it is just its own value.

\subsubsection*{Recurrence Relation}
For row $i$ ($2 \le i \le n$) and position $j$ ($1 \le j \le i$):

\[
    \text{dp}[i][j] = T[i][j] + \begin{cases}
        \text{dp}[i-1][1] & \text{if } j = 1 \\[6pt]
        \text{dp}[i-1][j-1] & \text{if } j = i \\[6pt]
        \max\!\big(\text{dp}[i-1][j-1],\; \text{dp}[i-1][j]\big) & \text{if } 1 < j < i
    \end{cases}
\]

\textbf{Justification:} To arrive at position $j$ in row $i$, we must have come from position $j-1$ or position $j$ in row $i-1$ (since a parent at position $k$ can move to $k$ or $k+1$). Therefore:
\begin{itemize}[nosep]
    \item If $j = 1$: the only parent is position $1$ in row $i-1$.
    \item If $j = i$: the only parent is position $j-1 = i-1$ in row $i-1$.
    \item Otherwise: we take the maximum of the two possible parents.
\end{itemize}

\subsubsection*{Optimal Value}
The answer is:
\[
    \text{Maximum Path Sum} = \max_{1 \le j \le n} \text{dp}[n][j]
\]
We scan the last row of the DP table and return the largest value.

\subsection*{Algorithm}

\begin{algorithm}[H]
\caption{\textsc{TriangleMaxPathSum}$(T, n)$}
\begin{algorithmic}[1]
\Require Triangle $T[1..n][1..i]$ and number of rows $n$
\Ensure Maximum path sum from top to bottom
\State $\text{dp}[1][1] \gets T[1][1]$
\For{$i \gets 2$ \textbf{to} $n$}
    \For{$j \gets 1$ \textbf{to} $i$}
        \If{$j = 1$}
            \State $\text{dp}[i][j] \gets T[i][j] + \text{dp}[i-1][1]$
        \ElsIf{$j = i$}
            \State $\text{dp}[i][j] \gets T[i][j] + \text{dp}[i-1][j-1]$
        \Else
            \State $\text{dp}[i][j] \gets T[i][j] + \max(\text{dp}[i-1][j-1],\; \text{dp}[i-1][j])$
        \EndIf
    \EndFor
\EndFor
\State \Return $\max_{1 \le j \le n} \text{dp}[n][j]$
\end{algorithmic}
\end{algorithm}

\subsection*{Complexity Analysis}

The algorithm has \textbf{time complexity} $O(n^2)$ since it fills an $n \times n$ triangular DP table with constant-time operations per cell. The \textbf{space complexity} is $O(n^2)$ for storing the DP table, though it can be optimized to $O(n)$ by using only two rows at a time since each row depends only on the previous row.

\newpage

\section*{Question 2: Coin Change Problem}

\subsection*{Problem Statement}
Given a set of coin denominations $C = \{c_0, c_1, \ldots, c_{n-1}\}$ where each $c_i$ is a positive integer, and assuming each coin denomination can be used an unlimited number of times, determine the minimum number of coins required to make a target value $V$ exactly. If no combination of coins can form the target value $V$, report that no solution exists.

\subsection*{Dynamic Programming Formulation}

\subsubsection*{Notation}
Let $C[i]$ denote the value of the $i$-th coin denomination (0-indexed), where $0 \le i < n$.

\subsubsection*{Subproblem Definition}
Define $\text{dp}[v]$ as the \textbf{minimum number of coins} required to make the value $v$, for $0 \le v \le V$.

\subsubsection*{Base Case}
\[
    \text{dp}[0] = 0
\]
No coins are needed to make a value of 0.

\subsubsection*{Recurrence Relation}
For each value $v$ where $1 \le v \le V$:

\[
    \text{dp}[v] = \begin{cases}
        \infty & \text{if no coin denomination can form } v \\[6pt]
        1 + \min_{i: C[i] \le v} \text{dp}[v - C[i]] & \text{otherwise}
    \end{cases}
\]

\textbf{Justification:} To form value $v$, we can use any coin $C[i]$ such that $C[i] \le v$, and then recursively form the remaining value $v - C[i]$. We choose the coin that results in the minimum total number of coins.

\subsubsection*{Optimal Value}
The answer is:
\[
    \text{Minimum Coins} = \begin{cases}
        \text{dp}[V] & \text{if } \text{dp}[V] < \infty \\
        \text{No solution exists} & \text{otherwise}
    \end{cases}
\]

\subsection*{Algorithm}

\begin{verbatim}
def min_coins(C, V):
    """
    Find the minimum number of coins required to make value V.
    """
    # dp[v] = minimum coins needed to make value v
    dp = [float('inf')] * (V + 1)
    dp[0] = 0  # Base case: 0 coins for value 0

    # Fill DP table for each value from 1 to V
    for v in range(1, V + 1):
        # Try using each coin denomination
        for coin in C:
            if coin <= v and dp[v - coin] != float('inf'):
                dp[v] = min(dp[v], 1 + dp[v - coin])

    # Return result
    return dp[V] if dp[V] != float('inf') else -1
\end{verbatim}

\subsection*{Complexity Analysis}

The algorithm has \textbf{time complexity} $O(V \cdot n)$ where $V$ is the target value and $n$ is the number of coin denominations. The outer loop iterates $V$ times, and for each value, we examine all $n$ coin denominations. The \textbf{space complexity} is $O(V)$ for storing the DP table that tracks the minimum coins needed for each value from 0 to $V$.

\newpage

\section*{Question 3: Matrix Chain Multiplication}

\subsection*{Problem Statement}
Given five matrices with the following dimensions:
\[
    M_1: 1 \times 5, \quad M_2: 5 \times 2, \quad M_3: 2 \times 3, \quad M_4: 3 \times 9, \quad M_5: 9 \times 4
\]
Find the minimum number of scalar multiplications required to compute $\prod_{j=1}^{5} M_j$ and an optimal parenthesization.

\subsection*{Dynamic Programming Formulation}

Let $p = [1, 5, 2, 3, 9, 4]$ be the dimension array so that matrix $M_i$ has dimensions $p[i] \times p[i+1]$ (0-indexed). Define $m[i][j]$ as the minimum number of scalar multiplications needed to compute $M_i \cdot M_{i+1} \cdots M_j$. The base case is $m[i][i] = 0$ (a single matrix needs no multiplications), and the recurrence for $i < j$ is:
\[
    m[i][j] = \min_{i \le k < j} \Big\{ m[i][k] + m[k+1][j] + p[i] \cdot p[k+1] \cdot p[j+1] \Big\}
\]
We also store $s[i][j] = k$ (the optimal split point) to reconstruct the parenthesization.

\subsection*{Algorithm}

\begin{verbatim}
def optimal_matrix_mult(matrices):
    N = len(matrices)
    m = [[float('inf')] * N for _ in range(N)]
    s = [[0] * N for _ in range(N)]

    for i in range(N):
        m[i][i] = 0

    for L in range(2, N + 1):          # chain length
        for i in range(N - L + 1):
            j = i + L - 1
            for k in range(i, j):
                cost = (m[i][k] + m[k + 1][j]
                        + matrices[i][0] * matrices[k][1] * matrices[j][1])
                if cost < m[i][j]:
                    m[i][j] = cost
                    s[i][j] = k

    def construct_parens(i, j):
        if i == j:
            return f"M{i + 1}"
        k = s[i][j]
        left = construct_parens(i, k)
        right = construct_parens(k + 1, j)
        return f"({left} {right})"

    return m[0][N - 1], construct_parens(0, N - 1), m, s
\end{verbatim}

\subsection*{Computation}

We fill the $m$ and $s$ tables by increasing chain length $L$:

\medskip
\noindent\textbf{Chain length $L = 2$:}
\begin{align*}
    m[0][1] &= p[0] \cdot p[1] \cdot p[2] = 1 \cdot 5 \cdot 2 = 10, \quad s[0][1] = 0 \\
    m[1][2] &= p[1] \cdot p[2] \cdot p[3] = 5 \cdot 2 \cdot 3 = 30, \quad s[1][2] = 1 \\
    m[2][3] &= p[2] \cdot p[3] \cdot p[4] = 2 \cdot 3 \cdot 9 = 54, \quad s[2][3] = 2 \\
    m[3][4] &= p[3] \cdot p[4] \cdot p[5] = 3 \cdot 9 \cdot 4 = 108, \quad s[3][4] = 3
\end{align*}

\noindent\textbf{Chain length $L = 3$:}
\begin{align*}
    m[0][2]: \quad k=0&: m[0][0] + m[1][2] + 1 \cdot 5 \cdot 3 = 0 + 30 + 15 = 45 \\
                  k=1&: m[0][1] + m[2][2] + 1 \cdot 2 \cdot 3 = 10 + 0 + 6 = \mathbf{16} \quad \Rightarrow s[0][2] = 1 \\[4pt]
    m[1][3]: \quad k=1&: m[1][1] + m[2][3] + 5 \cdot 2 \cdot 9 = 0 + 54 + 90 = \mathbf{144} \\
                  k=2&: m[1][2] + m[3][3] + 5 \cdot 3 \cdot 9 = 30 + 0 + 135 = 165 \quad \Rightarrow s[1][3] = 1 \\[4pt]
    m[2][4]: \quad k=2&: m[2][2] + m[3][4] + 2 \cdot 3 \cdot 4 = 0 + 108 + 24 = 132 \\
                  k=3&: m[2][3] + m[4][4] + 2 \cdot 9 \cdot 4 = 54 + 0 + 72 = \mathbf{126} \quad \Rightarrow s[2][4] = 3
\end{align*}

\noindent\textbf{Chain length $L = 4$:}
\begin{align*}
    m[0][3]: \quad k=0&: m[0][0] + m[1][3] + 1 \cdot 5 \cdot 9 = 0 + 144 + 45 = 189 \\
                  k=1&: m[0][1] + m[2][3] + 1 \cdot 2 \cdot 9 = 10 + 54 + 18 = 82 \\
                  k=2&: m[0][2] + m[3][3] + 1 \cdot 3 \cdot 9 = 16 + 0 + 27 = \mathbf{43} \quad \Rightarrow s[0][3] = 2 \\[4pt]
    m[1][4]: \quad k=1&: m[1][1] + m[2][4] + 5 \cdot 2 \cdot 4 = 0 + 126 + 40 = \mathbf{166} \\
                  k=2&: m[1][2] + m[3][4] + 5 \cdot 3 \cdot 4 = 30 + 108 + 60 = 198 \\
                  k=3&: m[1][3] + m[4][4] + 5 \cdot 9 \cdot 4 = 144 + 0 + 180 = 324 \quad \Rightarrow s[1][4] = 1
\end{align*}

\noindent\textbf{Chain length $L = 5$:}
\begin{align*}
    m[0][4]: \quad k=0&: m[0][0] + m[1][4] + 1 \cdot 5 \cdot 4 = 0 + 166 + 20 = 186 \\
                  k=1&: m[0][1] + m[2][4] + 1 \cdot 2 \cdot 4 = 10 + 126 + 8 = 144 \\
                  k=2&: m[0][2] + m[3][4] + 1 \cdot 3 \cdot 4 = 16 + 108 + 12 = 136 \\
                  k=3&: m[0][3] + m[4][4] + 1 \cdot 9 \cdot 4 = 43 + 0 + 36 = \mathbf{79} \quad \Rightarrow s[0][4] = 3
\end{align*}

\subsection*{DP Tables}

\noindent\textbf{Table $m$ (minimum costs):}
\begin{center}
\begin{tabular}{c|ccccc}
    $m[i][j]$ & $j=0$ & $j=1$ & $j=2$ & $j=3$ & $j=4$ \\ \hline
    $i=0$ & 0 & 10 & 16 & 43 & \textbf{79} \\
    $i=1$ & & 0 & 30 & 144 & 166 \\
    $i=2$ & & & 0 & 54 & 126 \\
    $i=3$ & & & & 0 & 108 \\
    $i=4$ & & & & & 0 \\
\end{tabular}
\end{center}

\noindent\textbf{Table $s$ (split points):}
\begin{center}
\begin{tabular}{c|ccccc}
    $s[i][j]$ & $j=0$ & $j=1$ & $j=2$ & $j=3$ & $j=4$ \\ \hline
    $i=0$ & -- & 0 & 1 & 2 & 3 \\
    $i=1$ & & -- & 1 & 1 & 1 \\
    $i=2$ & & & -- & 2 & 3 \\
    $i=3$ & & & & -- & 3 \\
    $i=4$ & & & & & -- \\
\end{tabular}
\end{center}

\subsection*{Result}

\[
    \text{Minimum scalar multiplications} = m[0][4] = \boxed{79}
\]

\noindent\textbf{Optimal Parenthesization:} Tracing back through the $s$ table:
\begin{itemize}[nosep]
    \item $s[0][4] = 3$: split into $(M_1 \cdots M_4)$ and $(M_5)$
    \item $s[0][3] = 2$: split into $(M_1 \cdots M_3)$ and $(M_4)$
    \item $s[0][2] = 1$: split into $(M_1 \cdot M_2)$ and $(M_3)$
\end{itemize}
\[
    \textbf{Optimal Parenthesization:} \quad ((((M_1 \, M_2) \, M_3) \, M_4) \, M_5)
\]

\subsection*{Complexity Analysis}

The algorithm has \textbf{time complexity} $O(n^3)$ due to three nested loops (chain length, start index, split point), where $n$ is the number of matrices. The \textbf{space complexity} is $O(n^2)$ for storing the $m$ and $s$ tables.

\newpage

\section*{Question 4: Edit Distance}

\subsection*{Problem Statement}
Find the edit distance between $s_1 = \texttt{INTENTION}$ and $s_2 = \texttt{EXECUTION}$ using dynamic programming. The allowed operations (each costing 1) are: \textbf{insert}, \textbf{delete}, and \textbf{substitute}.

\subsection*{Dynamic Programming Formulation}

Let $\text{dp}[i][j]$ be the edit distance between the first $i$ characters of $s_1$ and the first $j$ characters of $s_2$ (0-indexed). Base cases: $\text{dp}[i][0] = i$ (delete all) and $\text{dp}[0][j] = j$ (insert all). The recurrence for $i, j \ge 1$ is:
\[
    \text{dp}[i][j] = \begin{cases}
        \text{dp}[i-1][j-1] & \text{if } s_1[i-1] = s_2[j-1] \\[4pt]
        1 + \min\!\big(\text{dp}[i-1][j],\; \text{dp}[i][j-1],\; \text{dp}[i-1][j-1]\big) & \text{otherwise}
    \end{cases}
\]

\subsection*{DP Table}

With $s_1 = \texttt{INTENTION}$ (rows) and $s_2 = \texttt{EXECUTION}$ (columns):

\begin{center}
\small
\begin{tabular}{c|cccccccccc}
    & $\varepsilon$ & E & X & E & C & U & T & I & O & N \\ \hline
    $\varepsilon$ & 0 & 1 & 2 & 3 & 4 & 5 & 6 & 7 & 8 & 9 \\
    I & 1 & 1 & 2 & 3 & 4 & 5 & 6 & 6 & 7 & 8 \\
    N & 2 & 2 & 2 & 3 & 4 & 5 & 6 & 7 & 7 & 7 \\
    T & 3 & 3 & 3 & 3 & 4 & 5 & 5 & 6 & 7 & 8 \\
    E & 4 & 3 & 4 & 3 & 4 & 5 & 6 & 6 & 7 & 8 \\
    N & 5 & 4 & 4 & 4 & 4 & 5 & 6 & 7 & 7 & 7 \\
    T & 6 & 5 & 5 & 5 & 5 & 5 & 5 & 6 & 7 & 8 \\
    I & 7 & 6 & 6 & 6 & 6 & 6 & 6 & 5 & 6 & 7 \\
    O & 8 & 7 & 7 & 7 & 7 & 7 & 7 & 6 & 5 & 6 \\
    N & 9 & 8 & 8 & 8 & 8 & 8 & 8 & 7 & 6 & \textbf{5} \\
\end{tabular}
\end{center}
\subsection*{Algorithm}

\begin{verbatim}
def edit_distance(s1, s2):
    m, n = len(s1), len(s2)
    dp = [[0] * (n + 1) for _ in range(m + 1)]
    for i in range(m + 1): dp[i][0] = i
    for j in range(n + 1): dp[0][j] = j
    for i in range(1, m + 1):
        for j in range(1, n + 1):
            if s1[i - 1] == s2[j - 1]:
                dp[i][j] = dp[i - 1][j - 1]
            else:
                dp[i][j] = 1 + min(dp[i - 1][j], dp[i][j - 1], dp[i - 1][j - 1])
    return dp[m][n]
\end{verbatim}

\subsection*{Result}

\[
    \text{Edit Distance}(\texttt{INTENTION},\; \texttt{EXECUTION}) = \boxed{5}
\]

\newpage

\section*{Question 5: Longest Palindromic Subsequence}

\subsection*{Problem Statement}
Given a sequence $S = \langle s_1, \ldots, s_n \rangle$ over some alphabet $\Sigma$, design a dynamic programming algorithm to return the length of a longest palindromic subsequence in $S$.

\subsection*{Dynamic Programming Formulation}

Define $\text{dp}[i][j]$ as the length of the longest palindromic subsequence in the substring $S[i \ldots j]$ (0-indexed). Base cases: $\text{dp}[i][i] = 1$ for all $i$ (single characters are palindromes), and $\text{dp}[i][j] = 0$ when $i > j$. The recurrence for $i < j$ is:
\[
    \text{dp}[i][j] = \begin{cases}
        2 + \text{dp}[i+1][j-1] & \text{if } S[i] = S[j] \\[4pt]
        \max\!\big(\text{dp}[i+1][j],\; \text{dp}[i][j-1]\big) & \text{otherwise}
    \end{cases}
\]
The answer is $\text{dp}[0][n-1]$.

\subsection*{Algorithm}

\begin{verbatim}
def longest_palindromic_subseq(S):
    n = len(S)
    dp = [[0] * n for _ in range(n)]

    for i in range(n):
        dp[i][i] = 1

    for length in range(2, n + 1):
        for i in range(n - length + 1):
            j = i + length - 1
            if S[i] == S[j]:
                dp[i][j] = 2 + dp[i + 1][j - 1]
            else:
                dp[i][j] = max(dp[i + 1][j], dp[i][j - 1])

    return dp[0][n - 1]
\end{verbatim}

\subsection*{Complexity Analysis}

\textbf{Time Complexity:} The algorithm uses two nested loops. The outer loop iterates over subsequence lengths $\ell$ from 2 to $n$. For each length $\ell$, the inner loop iterates over all valid starting indices $i$, running $n - \ell + 1$ times. Each cell requires only constant-time work (one comparison and at most two table lookups). The total number of cells filled is $\frac{n(n-1)}{2} = O(n^2)$, so the overall time complexity is $O(n^2)$.

\textbf{Space Complexity:} The algorithm stores an $n \times n$ DP table, so the space complexity is $O(n^2)$.

\newpage

\section*{Question 6: Shortest Common Supersequence}

\subsection*{Problem Statement}
Given two sequences $x = \langle x_1, \ldots, x_m \rangle$ and $y = \langle y_1, \ldots, y_n \rangle$, find the length of a shortest common supersequence of $x$ and $y$---a shortest sequence $z$ such that both $x$ and $y$ are subsequences of $z$.

\subsection*{Dynamic Programming Formulation}

Define $\text{dp}[i][j]$ as the length of the shortest common supersequence of $x[0 \ldots i-1]$ and $y[0 \ldots j-1]$. Base cases: $\text{dp}[i][0] = i$ and $\text{dp}[0][j] = j$. The recurrence for $i, j \ge 1$ is:
\[
    \text{dp}[i][j] = \begin{cases}
        1 + \text{dp}[i-1][j-1] & \text{if } x[i-1] = y[j-1] \\[4pt]
        1 + \min\!\big(\text{dp}[i-1][j],\; \text{dp}[i][j-1]\big) & \text{otherwise}
    \end{cases}
\]
The answer is $\text{dp}[m][n]$.

\subsection*{Algorithm}

\begin{verbatim}
def shortest_common_superseq(x, y):
    m, n = len(x), len(y)
    dp = [[0] * (n + 1) for _ in range(m + 1)]

    for i in range(m + 1):
        dp[i][0] = i
    for j in range(n + 1):
        dp[0][j] = j

    for i in range(1, m + 1):
        for j in range(1, n + 1):
            if x[i - 1] == y[j - 1]:
                dp[i][j] = 1 + dp[i - 1][j - 1]
            else:
                dp[i][j] = 1 + min(dp[i - 1][j], dp[i][j - 1])

    return dp[m][n]
\end{verbatim}

\subsection*{Complexity Analysis}

\textbf{Time Complexity:} The algorithm uses two nested loops: the outer loop runs $m$ iterations and the inner loop runs $n$ iterations. Each cell requires only constant-time work (one comparison and at most two table lookups). The total number of cells filled is $m \cdot n$, so the overall time complexity is $O(mn)$.

\textbf{Space Complexity:} The algorithm stores an $(m+1) \times (n+1)$ DP table, so the space complexity is $O(mn)$.

\newpage

\section*{Question 7: Beautiful Sequence}

\subsection*{Problem Statement}
For a sequence $X = \langle x_1, x_2, \ldots, x_n \rangle$, the mirror of $X$ is $Y = \langle x_n, x_{n-1}, \ldots, x_1 \rangle$. A sequence is \textit{beautiful} if it equals its mirror. Design an $O(n^2)$ DP algorithm to find the minimum number of elements to add to $X$ (anywhere in the sequence) to make it beautiful.

\subsection*{Dynamic Programming Formulation}

The minimum insertions to make $X$ a palindrome equals $n - \text{LPS}(X)$, where $\text{LPS}(X)$ is the length of the longest palindromic subsequence. The LPS already forms a palindromic ``skeleton''; each remaining character needs exactly one insertion to mirror it.

Define $\text{dp}[i][j]$ as the length of the longest palindromic subsequence in $X[i \ldots j]$ (0-indexed). Base cases: $\text{dp}[i][i] = 1$, $\text{dp}[i][j] = 0$ when $i > j$. The recurrence for $i < j$ is:
\[
    \text{dp}[i][j] = \begin{cases}
        2 + \text{dp}[i+1][j-1] & \text{if } X[i] = X[j] \\[4pt]
        \max\!\big(\text{dp}[i+1][j],\; \text{dp}[i][j-1]\big) & \text{otherwise}
    \end{cases}
\]
The answer is $n - \text{dp}[0][n-1]$.

\subsection*{Algorithm}

\begin{verbatim}
def min_insertions_beautiful(X):
    n = len(X)
    dp = [[0] * n for _ in range(n)]

    for i in range(n):
        dp[i][i] = 1

    for length in range(2, n + 1):
        for i in range(n - length + 1):
            j = i + length - 1
            if X[i] == X[j]:
                dp[i][j] = 2 + dp[i + 1][j - 1]
            else:
                dp[i][j] = max(dp[i + 1][j], dp[i][j - 1])

    return n - dp[0][n - 1]
\end{verbatim}

\subsection*{Complexity Analysis}

The algorithm fills an $n \times n$ table, giving \textbf{time complexity} $O(n^2)$ and \textbf{space complexity} $O(n^2)$.

\newpage

\section*{Question 8: Subset Sum}

\subsection*{Problem Statement}
Given a set of nonnegative integers $S$ and a target integer $t$, design a dynamic programming algorithm to find a subset of $S$ whose sum equals $t$.

\subsection*{Dynamic Programming Formulation}

Let $S = \{s_0, s_1, \ldots, s_{n-1}\}$. Define $\text{dp}[i][j]$ as a boolean indicating whether it is possible to achieve sum $j$ using elements from $\{s_0, \ldots, s_{i-1}\}$. Base case: $\text{dp}[i][0] = \texttt{True}$ for all $i$ (empty subset sums to 0). The recurrence for $i \ge 1$ is:
\[
    \text{dp}[i][j] = \begin{cases}
        \text{dp}[i-1][j] \;\lor\; \text{dp}[i-1][j - s_{i-1}] & \text{if } s_{i-1} \le j \\[4pt]
        \text{dp}[i-1][j] & \text{otherwise}
    \end{cases}
\]
A solution exists if and only if $\text{dp}[n][t] = \texttt{True}$.

\subsection*{Algorithm}

\begin{verbatim}
def subset_sum(S, t):
    n = len(S)
    dp = [[False] * (t + 1) for _ in range(n + 1)]

    for i in range(n + 1):
        dp[i][0] = True

    for i in range(1, n + 1):
        for j in range(1, t + 1):
            dp[i][j] = dp[i - 1][j]
            if S[i - 1] <= j:
                dp[i][j] = dp[i][j] or dp[i - 1][j - S[i - 1]]

    if not dp[n][t]:
        return None

    subset = []
    j = t
    for i in range(n, 0, -1):
        if dp[i][j] and not dp[i - 1][j]:
            subset.append(S[i - 1])
            j -= S[i - 1]

    return subset
\end{verbatim}

\subsection*{Complexity Analysis}

\textbf{Time Complexity:} The algorithm uses two nested loops: the outer loop runs $n$ iterations (one per element) and the inner loop runs $t$ iterations (one per possible sum). Each cell requires only constant-time work (one comparison and a boolean OR). The total number of cells filled is $n \cdot t$, so the overall time complexity is $O(nt)$, where $n = |S|$ and $t$ is the target value.

\textbf{Space Complexity:} The algorithm stores an $(n+1) \times (t+1)$ DP table, so the space complexity is $O(nt)$.

\end{document}
